% Part: methods
% Chapter: induction
% Section: introduction

\documentclass[../../../include/open-logic-section]{subfiles}

\begin{document}

\olfileid{mth}{ind}{int}

\olsection{مقدمه}

استقرا فن مهمی در اثبات است که به شکل‌های گوناگون تقریباً در همهٔ
حوزه‌های منطق، علوم رایانهٔ نظری و ریاضیات به کار می‌رود. اثبات بسیاری
از نتایج منطق به آن نیاز دارد.

استقرا را اغلب در برابر استنتاج قرار می‌دهند و آن را استنتاج از خاص
به عام توصیف می‌کنند. برای نمونه، اگر زمردهای سبز بسیاری مشاهده کنیم
و هیچ چیزی که زمرد بنامیم و سبز نباشد نبینیم، شاید نتیجه بگیریم همهٔ
زمردها سبزند. این یک استنتاج استقرایی است، زیرا از موارد خاص متعدد
— این زمرد سبز است، آن زمرد سبز است و مانند آن — به گزاره‌ای کلی —
همهٔ زمردها سبزند — می‌رود. استقرای \emph{ریاضی} نیز استنتاجی است که
به گزاره‌ای کلی می‌رسد، اما ماهیت آن با این «استقرای ساده» بسیار
متفاوت است.

به‌طور بسیار تقریبی، یک اثبات استقرایی در ریاضیات نتیجه می‌گیرد که
همهٔ اشیای ریاضی از نوعی معین، خاصیتی معین دارند. در ساده‌ترین حالت،
اشیای ریاضی مورد نظر اعداد طبیعی‌اند. در آن حالت، اثبات استقرایی برای
نشان‌دادن این به کار می‌رود که همهٔ اعداد طبیعی خاصیتی دارند، و این
کار را با نشان‌دادن دو امر زیر انجام می‌دهد:
\begin{enumerate}
    \item عدد $0$ آن خاصیت را دارد، و
    \item هرگاه عدد~$k$ آن خاصیت را داشته باشد، $k+1$ نیز آن را دارد.
\end{enumerate}
استقرا بر اعداد طبیعی را اغلب می‌توان برای اثبات گزاره‌های کلی دربارهٔ
اشیای ریاضی‌ای نیز به کار برد که می‌توان عددی به آن‌ها نسبت داد. برای
نمونه، هر مجموعهٔ متناهی تعداد متناهی~$n$ عضو دارد، و اگر بتوانیم با
استقرا نشان دهیم هر عدد~$n$ دارای این خاصیت است که «همهٔ مجموعه‌های
متناهی با اندازهٔ~$n$، \dots هستند»، آنگاه چیزی دربارهٔ همهٔ مجموعه‌های
متناهی نشان داده‌ایم.

استقرا را می‌توان به اشیای ریاضی‌ای نیز تعمیم داد که
\emph{به‌طور استقرایی تعریف شده‌اند}. برای نمونه، عبارت‌های یک زبان
صوری، مانند عبارت‌های منطق مرتبهٔ اول، به‌طور استقرایی تعریف می‌شوند.
\emph{استقرای ساختاری} راهی برای اثبات نتایج دربارهٔ همهٔ چنین عبارت‌هایی
است. استقرای ساختاری به‌ویژه در منطق بسیار سودمند و به‌طور گسترده
به‌کاررفته است.

\end{document}
